%% 논문 메타 정보 — 아래 값들을 확인/수정하세요.
%% (TODO 표시는 학과 사무실·지도교수와 확인이 필요한 항목)

% 제목 (국문 / 영문). 줄바꿈이 필요하면 \\ 사용
\pnutitle{분 단위 공개 텔레메트리 기반\\리그 오브 레전드 교전 결과 예측:\\킬 조건부 교전 국소화와 다중 모달 학습 프레임워크}
\pnutitleen{Engagement Outcome Prediction in League of Legends\\from Minute-Resolution Public Telemetry:\\Kill-Conditioned Localization and a Multi-Modal Learning Framework}

% 저자
\pnuauthor{백성은}            % 붙여 쓴 이름 (인준지 "OOO의 ... 인준함")
\pnuauthorspaced{백\ 성\ 은}   % 띄어 쓴 이름 (표지)
\pnuauthoren{Seongeun Baek}

% 지도교수
\pnuadvisor{권\ 준\ 호}
\pnuadvisoren{Joonho Kwon}

% 학과 — TODO: 대학원 소속 학과(전공) 정식 명칭 확인
%   후보: 정보융합공학과 (컴퓨터공학전공 / AI전공)
\pnudept{정보융합공학과}
\pnudepten{Information Convergence Engineering}

% 학위명: "공학" → 공학석사
\pnudegree{공학}

% 학위수여년월 / 인준(최종심사) 연월일 — TODO: 실제 일정으로 수정
\pnugraddate{2027년\ 2월}
\pnuapprovaldate{2026년\ 12월\ \ \ 일}

% 심사위원 (위원장, 위원, 위원) — 비워 두면 밑줄만 출력
\pnucommittee{}{}{}

% 주요어 (국문 초록 끝) / Keywords (영문 초록 끝)
\pnukeywords{리그 오브 레전드, 교전 결과 예측, 이스포츠 분석, 시공간 그래프 신경망, 그래디언트 부스팅, 공개 텔레메트리}
\pnukeywordsen{League of Legends, engagement outcome prediction, esports analytics, spatio-temporal graph neural networks, gradient boosting, public telemetry}
